\section{Multiple sequence columns and sequence joins}
\label{sec:appendix-joins}

A sequence search has two roles: the first argument of \texttt{mmsearch} is
the target, and the second is the query. SQL aliases name those roles when
both arguments come from stored rows. A table can contain several sequence
columns; every search and sequence operation names the column it uses.
Section~\ref{sec:mmsearch-join-ownership} explains how the custom scan keeps
the resulting alignment fields with the correct search.

\begin{itemize}
  \item \textbf{One column and a literal.} Search one selected sequence
    column against a query written in SQL.
  \item \textbf{One table and two aliases.} Use one alias for targets and
    another for stored queries. This is a self-join.
  \item \textbf{Two tables.} Keep the targets and queries in separate
    relations, then join their sequence columns.
  \item \textbf{Several searches.} Give each search a supported place in
    the plan and an unambiguous target. Two predicates on the same target
    cannot generally share one set of alignment fields.
\end{itemize}

\begin{figure}[htbp]
\centering
\begin{tikzpicture}[font=\small,
  jt/.style={draw=appendixBlue,fill=appendixBlue!7,minimum width=2.2cm,
    minimum height=.65cm,align=center},
  jr/.style={draw=appendixOrange,fill=appendixOrange!10,minimum width=2.2cm,
    minimum height=.65cm,align=center},
  ja/.style={-{Latex[length=1.8mm]},thick}]
  \node[font=\bfseries] at (2.65,3.75) {Self-join};
  \node[jt,minimum width=4.5cm] (one) at (2.65,2.8)
    {one table: \texttt{paired\_reads}};
  \node[jr] (st) at (1.25,1.45) {\texttt{targets.seq0}};
  \node[jr] (sq) at (4.05,1.45) {\texttt{probes.seq1}};
  \draw[ja] (one.south west) -- (st.north);
  \draw[ja] (one.south east) -- (sq.north);
  \draw[ja] (sq.west) -- (st.east);
  \node[align=center,text=black!70] at (2.65,.35)
    {Two SQL roles can read\\different rows of the same table.};
  \draw[black!25] (5.75,0) -- (5.75,4.05);
  \node[font=\bfseries] at (8.85,3.75) {Two-table join};
  \node[jt] (twoa) at (7.45,2.8) {target table};
  \node[jt] (twob) at (10.25,2.8) {query table};
  \node[jr] (tt) at (7.45,1.45) {\texttt{t.sequence}};
  \node[jr] (tq) at (10.25,1.45) {\texttt{q.sequence}};
  \draw[ja] (twoa) -- (tt);
  \draw[ja] (twob) -- (tq);
  \draw[ja] (tq.west) -- (tt.east);
  \node[align=center,text=black!70] at (8.85,.35)
    {The same target/query roles\\come from different tables.};
\end{tikzpicture}
\caption{The horizontal arrows show the query being searched against the
target. A self-join changes the SQL roles; it does not create a second
physical table. Neither layout promises a particular number of storage
reads or a particular join plan.}
\label{fig:appendix-join-roles}
\end{figure}

\subsection{Prepare and inspect a table with two sequence columns}

The examples below use five rows. Both columns contain DNA, but
\texttt{seq0} holds longer targets and \texttt{seq1} holds shorter queries.
The two strings on adjacent lines within a single sequence value are joined
by SQL; they are wrapped only to keep the listing readable.

\par\noindent\begin{minipage}{\linewidth}
\begin{lstlisting}[style=bioseq-appendix-sql]
CREATE TABLE paired_reads (
  pk integer PRIMARY KEY,
  sample_id text NOT NULL,
  seq0 bioseq.seq('dna4', 32) NOT NULL,
  seq1 bioseq.seq('dna4', 32) NOT NULL
) USING bioseq;

SET bioseq.index.insert_maintenance.enabled = on;
SET bioseq.vectorization.enabled = on;

INSERT INTO paired_reads VALUES
  (1, 'alpha',
   'TCACTGTAGGATCGTAGCTAGTCCGTACTGTAGCTGATCGTAGCCCTAG'
   'TCACTGTAGCCTGATCGTAGCTAGTCACTGTAGCCATACCTGAA',
   'CTAGTCACTGTAGCCTGATCGTAGCTAGT'),
  (2, 'alpha',
   'CCCCAACTTAGTCACTTCTCAAAAATTACTCTCTCATCATCATCTACTCC'
   'TAGTCACTGTAGCCTGATCGTAGCTAGTCACACACACCCCAAA',
   'TCGTAGCTAGGA'),
  (3, 'beta',
   'GGGGGTCACTTCTCAAAAATTAGGGGGGGGGTAGTCACTGTAGCCTGAT'
   'CGTAGCTAGGACACATATCATCATCTACTCCTATCATCATCTACTCCTA',
   'TCATCATCTACTCCTA'),
  (4, 'beta',
   'TTTCATCATCATCATCTACTCCTATCTATCATCATCATCTACTCCTATC'
   'TACTCTCATCATCCTGATCGTAGCTATACTCCTACTATCCTATT',
   'TGTAGCCTGATCGTA'),
  (5, 'gamma',
   'ACGTCTAGTCACTGTAGCCTGATCGTAGCTAGTGTACTCTAGTCACTGT'
   'AGCCTGATCGTAGCTAGTTGATCCCCATATCTTCACTAGTCACTGTAGC'
   'CTGATCGTAGCTAGTTTTTTTA',
   'CACTGTAGCCTGATCGTA');

SELECT pk, bioseq.len(seq0), bioseq.len(seq1)
FROM paired_reads ORDER BY pk;

SELECT pk, sample_id, seq1::text
FROM paired_reads WHERE sample_id = 'alpha' ORDER BY pk;

SELECT pk, seq0[:12]::text AS target_prefix,
       seq1[-8:]::text AS query_suffix
FROM paired_reads ORDER BY pk;
\end{lstlisting}
\end{minipage}

Each expression selects its own column. The last query returns up to twelve
letters from the start of \texttt{seq0} and eight from the end of
\texttt{seq1}. It does not perform alignment.

The two settings control different work:

\begin{itemize}
  \item \texttt{bioseq.index.insert\_maintenance.enabled} permits eligible
    index maintenance during insertion, subject to the index's update
    policy. It does not build every index or make all index schemes support
    per-row updates.
  \item \texttt{bioseq.vectorization.enabled} allows vectorized execution
    where the selected pipeline supports it. It does not change which
    column or alias is the target.
\end{itemize}

The configuration reference documents their defaults and limits.

\subsection{Choose the searched column and return its alignment fields}

These are separate searches. Changing the first argument changes the target
column; there is no implicit ``current sequence'' column.

\par\noindent\begin{minipage}{\linewidth}
\begin{lstlisting}[style=bioseq-appendix-sql]
SELECT pk FROM paired_reads
WHERE public.mmsearch(
  seq0, 'CTAGTCACTGTAGCCTGATCGTAGCTAGT'
          ::bioseq.seq('dna4', 32));

SELECT pk FROM paired_reads
WHERE public.mmsearch(
  seq1, 'TCATCATCTACTCCTA'::bioseq.seq('dna4', 32));
\end{lstlisting}
\end{minipage}

Project metadata and alignment fields in the same search query:

\par\noindent\begin{minipage}{\linewidth}
\begin{lstlisting}[style=bioseq-appendix-sql]
SELECT pk, sample_id,
       bioseq.target_start(pk) AS target_start,
       bioseq.target_end(pk) AS target_end,
       bioseq.query_start(pk) AS query_start,
       bioseq.query_end(pk) AS query_end,
       bioseq.cigar(pk) AS cigar
FROM paired_reads
WHERE public.mmsearch(
  seq0, 'CTAGTCACTGTAGCCTGATCGTAGCTAGT'
          ::bioseq.seq('dna4', 32),
  public.aligner_options(best_n => 3))
ORDER BY pk, target_start, query_start;
\end{lstlisting}
\end{minipage}

Here \texttt{pk} identifies a row from the only target alias, so its
alignment helpers have one possible search result. The key itself is not
aligned. A target can produce several result rows with the same key. Keep
the coordinates and other required alignment fields if those occurrences
matter. The options object overrides the corresponding GUC defaults for
this call; Section~\ref{sec:appendix-search-options} defines
\texttt{best\_n} and the other fields.

\subsection{Use a self-join for stored queries}

The aliases below give the table two roles. For each candidate pair,
\texttt{targets.seq0} is the target and \texttt{probes.seq1} is the query.

\par\noindent\begin{minipage}{\linewidth}
\begin{lstlisting}[style=bioseq-appendix-sql]
SELECT targets.pk AS target_pk, probes.pk AS probe_pk,
       bioseq.target_start(targets.pk) AS target_start,
       bioseq.target_end(targets.pk) AS target_end,
       bioseq.cigar(targets.pk) AS cigar
FROM paired_reads AS targets
CROSS JOIN paired_reads AS probes
WHERE public.mmsearch(targets.seq0, probes.seq1)
ORDER BY target_pk, probe_pk, target_start;
\end{lstlisting}
\end{minipage}

Remove the three helper columns when only the key pairs are needed. The
search can still return several alignments for one pair. Use
\texttt{SELECT DISTINCT} on the two keys only when the required result is
one row per matching pair, rather than one row per alignment.

An explicit result stage makes later metadata filters easy to read:

\par\noindent\begin{minipage}{\linewidth}
\begin{lstlisting}[style=bioseq-appendix-sql]
WITH raw_matches AS MATERIALIZED (
  SELECT targets.pk AS target_pk, probes.pk AS probe_pk,
         targets.sample_id AS target_sample,
         probes.sample_id AS probe_sample
  FROM paired_reads AS targets
  CROSS JOIN paired_reads AS probes
  WHERE public.mmsearch(targets.seq0, probes.seq1)
)
SELECT * FROM raw_matches
WHERE target_sample <> probe_sample AND probe_pk IN (1, 3, 5)
ORDER BY target_pk, probe_pk;
\end{lstlisting}
\end{minipage}

\begin{itemize}
  \item \textbf{Filter results.} The outer conditions select already
    generated matches. This example does not avoid alignment for the rows
    discarded by those conditions.
  \item \textbf{Reduce input work.} If fewer input rows should reach the
    search, inspect the actual child plans. A clear alternative is to load
    the selected inputs into separate BioSeqDB tables first.
  \item \textbf{Do not assume all subqueries are new search targets.}
    A materialized result containing sequence values does not automatically
    acquire the storage and search behavior of a BioSeqDB relation.
\end{itemize}

\subsection{Remove pairs with the same row key}

To compare different rows, remove pairs whose keys are equal. The
\emph{diagonal} means precisely those equal-key pairs in this self-join.

\par\noindent\begin{minipage}{\linewidth}
\begin{lstlisting}[style=bioseq-appendix-sql]
WITH pairs AS MATERIALIZED (
  SELECT targets.pk AS target_pk, probes.pk AS probe_pk
  FROM paired_reads AS targets
  CROSS JOIN paired_reads AS probes
  WHERE public.mmsearch(targets.seq0, probes.seq1)
)
SELECT * FROM pairs
WHERE target_pk <> probe_pk
ORDER BY target_pk, probe_pk;
\end{lstlisting}
\end{minipage}

If the application defines a symmetric relationship, a condition such as
\texttt{target\_pk < probe\_pk} can retain one direction of each pair.
That assumption does not follow from this example: the target and query
use different columns and the alignment roles are directional. Removing
one direction without that assumption can discard distinct results.

\subsection{Put targets and queries in separate tables}

The same search roles can come from two physical tables. This example
copies the previously loaded values to keep the complete workflow runnable.

\par\noindent\begin{minipage}{\linewidth}
\begin{lstlisting}[style=bioseq-appendix-sql]
CREATE TABLE search_targets (
  target_id integer PRIMARY KEY,
  cohort text NOT NULL,
  sequence bioseq.seq('dna4', 32) NOT NULL
) USING bioseq;

CREATE TABLE search_queries (
  query_id integer PRIMARY KEY,
  query_name text NOT NULL,
  sequence bioseq.seq('dna4', 32) NOT NULL
) USING bioseq;

INSERT INTO search_targets
SELECT pk, sample_id, seq0 FROM paired_reads;
INSERT INTO search_queries
SELECT pk, sample_id, seq1 FROM paired_reads;

SELECT q.query_id, q.query_name, t.target_id, t.cohort,
       bioseq.cigar(t.target_id) AS cigar
FROM search_queries AS q
JOIN search_targets AS t
  ON public.mmsearch(t.sequence, q.sequence)
ORDER BY q.query_id, t.target_id;
\end{lstlisting}
\end{minipage}

For a supported inner-join shape, the planner can select
\texttt{BioseqBatchJoin}. Other supported shapes use a parameterized
\texttt{BioseqScan}: its query sequence comes from the current outer row.
SQL spelling alone does not guarantee one of these plans.

\subsection{Read the join plan and its capacity information}

\par\noindent\begin{minipage}{\linewidth}
\begin{lstlisting}[style=bioseq-appendix-sql]
EXPLAIN (VERBOSE, COSTS OFF)
SELECT q.query_id, t.target_id
FROM search_queries AS q
JOIN search_targets AS t
  ON public.mmsearch(t.sequence, q.sequence);
\end{lstlisting}
\end{minipage}

Add \texttt{ANALYZE} to execute the query and obtain actual counts. The
following fields are useful when the selected node is
\texttt{BioseqBatchJoin}; labels are abbreviated by omitting their common
\texttt{Bioseq} prefix.

\begin{center}
\small
\begin{tabular}{@{}p{.35\linewidth}p{.60\linewidth}@{}}
\toprule
Plan field & Meaning \\
\midrule
Query Batch Capacity & Maximum query count in one execution batch. It is
  not a limit on all retained query rows. \\
Residual Mmsearch Qual & \texttt{absent} records that the owned search
  marker was removed from ordinary Boolean filtering. \\
Index Read Requirement & The selected rule for using, requiring, or
  avoiding persisted target indexes. \\
Query Planner Width & PostgreSQL's estimate of bytes per query-side row. \\
Distinct Targets & Target child rows with a non-null sequence. Repeated
  target identities count again. \\
Query Batches & Number of query batches processed during execution. \\
\bottomrule
\end{tabular}
\end{center}

The join retains its query-side rows before processing the target side.
Therefore, a small batch capacity does not imply a small total memory
requirement. The plan also reports estimated admission bytes and actual
query materialization bytes. These help distinguish memory retained for
all queries from work grouped into one batch.

Check the custom node as well as its children. An ordinary executable
\texttt{Join Filter: mmsearch(...)} is not the intended replacement for an
alignment-owning custom node. Other pair conditions can remain as ordinary
filters and be evaluated after an alignment has been generated.

\subsection{Use LATERAL to express a query-dependent search}

\texttt{LATERAL} allows the inner query to refer to a preceding input.
Project the alignment fields inside the search query, then return those
ordinary output columns outside it:

\par\noindent\begin{minipage}{\linewidth}
\begin{lstlisting}[style=bioseq-appendix-sql]
SELECT q.query_id, matches.target_id, matches.cigar
FROM search_queries AS q
CROSS JOIN LATERAL (
  SELECT t.target_id, bioseq.cigar(t.target_id) AS cigar
  FROM search_targets AS t
  WHERE public.mmsearch(t.sequence, q.sequence)
) AS matches
ORDER BY q.query_id, matches.target_id;
\end{lstlisting}
\end{minipage}

A target metadata condition, for example \texttt{t.cohort = 'alpha'}, can
be added inside that subquery. Check the plan to see where it is evaluated
and whether it reduces the target input before alignment. PostgreSQL can
flatten a simple lateral subquery, so \texttt{LATERAL} alone does not force
a separate physical search stage.

A global \texttt{LIMIT} limits the whole result. It does not request a
fixed number of alignments for each query sequence. If the application
needs a per-query limit, specify its ordering and stage deliberately, then
check the plan and returned groups; do not infer that behavior from the
batch size or from the presence of \texttt{LATERAL}.

\subsection{Chain searches with separate target aliases}

The first search below aligns each query against a middle target. The
second uses that middle sequence as its query and searches a new target
alias. The roles, and the two alignment results, are distinct.

\par\noindent\begin{minipage}{\linewidth}
\begin{lstlisting}[style=bioseq-appendix-sql]
CREATE TABLE second_search_targets (
  target_id integer PRIMARY KEY,
  cohort text NOT NULL,
  sequence bioseq.seq('dna4', 32) NOT NULL
) USING bioseq;

INSERT INTO second_search_targets
SELECT target_id, cohort, sequence FROM search_targets;

SELECT q.query_id, middle_targets.target_id AS middle_id,
       final_targets.target_id AS final_id
FROM search_queries AS q
JOIN search_targets AS middle_targets
  ON public.mmsearch(middle_targets.sequence, q.sequence)
JOIN second_search_targets AS final_targets
  ON public.mmsearch(final_targets.sequence,
                     middle_targets.sequence)
ORDER BY q.query_id, middle_id, final_id;
\end{lstlisting}
\end{minipage}

Inspect the same statement with \texttt{EXPLAIN (ANALYZE, VERBOSE)} to
find the owner of each search. A later search in this shape can use a
parameterized \texttt{BioseqScan}; do not expect two
\texttt{BioseqBatchJoin} nodes simply because the SQL contains two joins.
Read execution counts, index choices, and algorithm fields at each
search-owning node.

\subsection{Keep alignment projections unambiguous}

A key such as \texttt{t.target\_id} can identify the current alignment
only when one search result is associated with that target in its query
scope. A key does not choose between two alignments from different
searches. When two candidate searches share the same source, an ambiguous
helper binding is rejected with SQLSTATE \texttt{42702}.

Keep helpers in the query that owns their search. For example, this
lateral query explicitly retains a subquery boundary with \texttt{OFFSET 0}
and exports ordinary coordinate columns:

\par\noindent\begin{minipage}{\linewidth}
\begin{lstlisting}[style=bioseq-appendix-sql]
SELECT q.query_id, matches.target_id,
       matches.target_start, matches.target_end
FROM search_queries AS q
CROSS JOIN LATERAL (
  SELECT t.target_id,
         bioseq.target_start(t.target_id) AS target_start,
         bioseq.target_end(t.target_id) AS target_end
  FROM search_targets AS t
  WHERE public.mmsearch(t.sequence, q.sequence)
  OFFSET 0
) AS matches
ORDER BY q.query_id, matches.target_id, matches.target_start;
\end{lstlisting}
\end{minipage}

The outer query reads \texttt{matches.target\_start}; it does not try to
recover a search result by calling a helper on a detached identifier.
This is a useful boundary for result expressions, not a promise that all
possible lateral or nested search combinations are supported.

\subsection{Find rows that match two queries without stacking searches}

The following form is unsupported: two ordinary search predicates ask the
same target binding to supply independent match results.

\par\noindent\begin{minipage}{\linewidth}
\begin{lstlisting}[style=bioseq-appendix-sql]
-- Unsupported: two searches on the same target binding.
SELECT pk FROM paired_reads
WHERE public.mmsearch(
        seq0, 'CTAGTCACTGTAGCCTGATCGTAGCTAGT'
                ::bioseq.seq('dna4', 32))
  AND public.mmsearch(
        seq0, 'TCATCATCTACTCCTA'::bioseq.seq('dna4', 32));
\end{lstlisting}
\end{minipage}

The ordinary stacked form is rejected with SQLSTATE \texttt{0A000}.
Section~\ref{sec:mmsearch-join-ownership} shows distinct target aliases
for independent searches. When only the set of matching row keys matters,
another simple solution is to run one search per result table and intersect
their keys:

\par\noindent\begin{minipage}{\linewidth}
\begin{lstlisting}[style=bioseq-appendix-sql]
CREATE TEMP TABLE first_matches AS
SELECT DISTINCT pk FROM paired_reads
WHERE public.mmsearch(
  seq0, 'CTAGTCACTGTAGCCTGATCGTAGCTAGT'
          ::bioseq.seq('dna4', 32));

CREATE TEMP TABLE second_matches AS
SELECT DISTINCT pk FROM paired_reads
WHERE public.mmsearch(
  seq0, 'TCATCATCTACTCCTA'::bioseq.seq('dna4', 32));

SELECT a.pk
FROM first_matches AS a JOIN second_matches AS b USING (pk)
ORDER BY a.pk;
\end{lstlisting}
\end{minipage}

This deliberately discards multiple occurrences within each search. To
retain them, project the required coordinates and CIGAR in each search
stage and join those result rows. The join can then return several
combinations for one key. If concurrent changes matter across the separate
statements, run the workflow in a transaction with the required snapshot
semantics.

\subsection{Search slices and convert their coordinates}

A transformed target is a logical sequence. Its reported positions refer
to that sequence. For \texttt{targets.seq0[16:-16]}, position zero is the
first retained letter, sixteen letters into the original value. The end
bound excludes the final sixteen letters.

\par\noindent\begin{minipage}{\linewidth}
\begin{lstlisting}[style=bioseq-appendix-sql]
WITH relative_matches AS MATERIALIZED (
  SELECT targets.pk AS target_pk, probes.pk AS probe_pk,
         16 AS slice_start,
         bioseq.target_start(targets.seq0[16:-16]) AS rel_start,
         bioseq.target_end(targets.seq0[16:-16]) AS rel_end,
         bioseq.cigar(targets.seq0[16:-16]) AS cigar
  FROM paired_reads AS targets
  CROSS JOIN paired_reads AS probes
  WHERE public.mmsearch(targets.seq0[16:-16], probes.seq1)
)
SELECT target_pk, probe_pk,
       slice_start + rel_start AS original_start,
       slice_start + rel_end AS original_end, cigar
FROM relative_matches
ORDER BY target_pk, probe_pk, original_start;
\end{lstlisting}
\end{minipage}

The helper arguments repeat the exact target slice. The outer query adds
the known start offset to the already projected coordinates. This example
uses zero-based, half-open ranges: the start is included and the end is
excluded. For a negative or computed start bound, first resolve that bound
against the original sequence length; adding its unevaluated SQL spelling
would be wrong. Slicing can change index eligibility and the selected plan,
so inspect a transformed search separately from the unsliced version.

\subsection{Check the workload and remove the example tables}

Before timing a larger join, check the following:

\begin{enumerate}
  \item Count the query and target rows that actually reach the search.
    A filter on emitted matches is not necessarily an input restriction.
  \item Find the search-owning custom nodes. Check the query batch
    capacity and the memory used to retain all query-side rows.
  \item Use \texttt{EXPLAIN ANALYZE} to inspect actual batch counts,
    distinct targets, index source, and selected algorithms. It executes
    the statement, unlike plain \texttt{EXPLAIN}.
  \item Decide whether the output should contain all alignments or only
    distinct row pairs. \texttt{DISTINCT} changes the result and adds work.
  \item Distinguish a limit on the whole output from a limit per query.
    Do not interpret \texttt{LIMIT} as a bound on alignment work.
\end{enumerate}

After running the examples, remove their tables and restore the two session
defaults:

\par\noindent\begin{minipage}{\linewidth}
\begin{lstlisting}[style=bioseq-appendix-sql]
DROP TABLE IF EXISTS first_matches, second_matches;
DROP TABLE IF EXISTS second_search_targets;
DROP TABLE IF EXISTS search_queries, search_targets;
DROP TABLE IF EXISTS paired_reads;
RESET bioseq.index.insert_maintenance.enabled;
RESET bioseq.vectorization.enabled;
\end{lstlisting}
\end{minipage}
